\documentclass[../main.tex]{subfiles}

\graphicspath{{\subfix{../images/}}}



\begin{document}

%% BEGIN SUBFILE %%

\newpage
\section{Operational Amplifiers}

Black box approach:

\begin{enumerate}[label=\textbf{-}]
	\item 2 input inverting (-) and non-inverting (+)
	\item 2 Power supply terminals $V_{cc}$ or $V_{dd}$ and $V_{ee}$ or $V_{ss}$ 
	\item 1 output terminal
	\item $V_{out} = A_d  \cdot v_d$ $A_d$ is the open loop gain of the amplifer
\end{enumerate}

The amplifier has a linear region where this relation applies before reaching saturation implied by power supply voltages. Since the open loop gain is very large $A_d \gg 100'000$ and we usually supply $\pm$ 12 or 15 V we would get this linear region:

\begin{equation}
	\Delta v_d = \frac{15 - (-15)}{100'000} = 300 \mu V
\end{equation}

Very little swing of the differental voltage, that is the differnce between the inverting and non inverting terminal on the amplifier.


\textbf{Ideal Op-Amp}: let's now make some assumptions around the amplifier:

\begin{enumerate}[label=\textbf{-}]
	\item $A_d \rightarrow \infty$ but since $V_{out} = A_d \cdot v_d$
	\item If $A_d$ is infinite and the $V_{out}$ is finite in the linear region means that the differential voltage $v_d$ tents to 0.
	\item We also assume that the input absorb no current $i_-$ $i_+$ is 0.
	\item We also assume the $r_d$ is infinite and that the $r_0$ is zero. This means that the op-amp absorbs no current and can drive ideally any circuit.
 \end{enumerate}


Ideally we have the + and - terminal such as that the $V_{out}$ is in phase with. the voltage on + terminal. Also if V on + an - is equal ideally we have an output of 0, this is called common-mode rejection. This mean an ideal op-amp has a zero common mode gain or infinite common mode rejection. This makes the op amp a \textbf{differential-input single-ended-output}. The gain A is called open-loop or differetial gain.\\

The op-amp are dc coupled meaninf directly coupled. meaning it can amplify as low as freq at 0 Hz.\\

Regarding the bandwidth the gain A remains constant all the way. Ideally it amplifies from DC signal to infinite frequencies. Also the gain A (open loop) is ideally very large ideally infinite.

But how do we use it? We close the loop.


\subsection{Inverting configuration}

We use a network of passive componets to close the loop:
We put the voltage on the - terminal and feed back the output to it:


\vspace*{7cm}


We define now the closed loop gain as $G = \frac{V_o}{V_i}$
Ideally we will have:

\begin{equation}
	v_2 - v_1 = \frac{v_O}{A} = 0
\end{equation}

since A is supposed to be infinitly large.
ideally $v_1$ approached and areached $v_2$ "tracking each other in potential" and we speak of a "virtual short circuit" it means that what is on 2 will appear on 1. Since 2 is GND terminal one will be called a virtual ground.

We apply Ohm's law to the resistors:

\begin{equation}
	i_1 = \frac{v_i - v_1}{R_1} = \frac{V_i - 0}{R_1} = \frac{V_i}{R_1}
\end{equation}

sice the terminal has infinite impedance it will go via the $R_2$ resistance.

We have that:

\begin{equation}
	v_1 = 0 \rightarrow i_1 = \frac{v_I}{R_1} \rightarrow i_2 = i_1 = \frac{v_I}{R_1} \rightarrow v_o = 0 - \frac{v_I}{R_1} R_2 \rightarrow v_O = - \frac{R_2}{R_1} v_I
\end{equation}

Closed loop gain $\rightarrow$ ration of the resistor values. This is the inverting configuration.


We then define the gain $G = - \frac{R_2}{R_1}$ account for possible $R_s$.


\subsection{NON Inverting configuration}

We apply now thw voltage directly to the NON inverting input of the opamp while the inverting terminal is connected to GND via. $R_1$.\\


\vspace*{7cm}


We will have:

\begin{equation}
	v_+ = v_{in} \rightarrow v_d = v{+} - v_- \rightarrow v_- = v_+ - v_d = v_+ = v_{in} \rightarrow i_1 = \frac{V_{in}}{R_1}
\end{equation}

then:

\begin{equation}
	V_{r2} = i_1 \cdot R_2 = v_{in} \frac{R_2}{R_1}
\end{equation}

Since:

\begin{equation}
	V_{in} + v_{r2} - v_{out} = 0 \rightarrow V_{out} = v_{in} + v_{in} \frac{R_2}{R_1} = v_{in} (1 + \frac{R_2}{R_1})
\end{equation}


\subsection{CONTROL THEORY}

Design basic control blocks:

I is input and Y is ouput.
The relation:

\begin{equation}
	\frac{Y}{I} = \frac{1}{\beta} \cdot \frac{A\beta}{1 + A\beta}
\end{equation}

We remind that A is the open loop gain while the $\beta$ term is what is called the negative feedback factor.

\begin{equation}
	e = I - \beta Y \rightarrow Y = Ae \rightarrow Y = A(I - \beta Y)
\end{equation}

Expand move factor on one side and factor Y therms we obtain that:

\begin{equation}
	\frac{Y}{I} = \frac{1}{\beta} \cdot \frac{A\beta}{1 + A\beta} = \frac{A}{1 + A\beta}
\end{equation}

The $A\beta$ is the so-called loop gain.
Since the $A\beta \rightarrow \infty$ therfore:

\begin{equation}
	\frac{Y}{I} \approx \frac{1}{\beta}
\end{equation}

We find that $\beta$:

\begin{equation}
	\beta = \frac{R_1}{R_1 + R_2} \rightarrow \frac{1}{\beta} = \frac{R_1+ R_2}{R_1} = 1 + \frac{R_2}{R_1}
\end{equation}


\subsection{Voltage follower}

We now have a gain G of 1 if $R_2$ is 0 and $R_1 = \infty$



\subsection{NON IDEAL OP-AMP}

We now study the op amp when $A_d$ is not infinite menaing that $V_d$ is not 0.


\vspace*{10cm}


We start from the NON inverting voltage amplifier.








%% END SUBFILE %%

\end{document}